\chapter{Conclusions and future work}
\label{ch:conclusions}

\section{Conclusions}
\label{sec:concl-summary}

This work implemented two top-k selection algorithms on three architectures,
using a common harness for their evaluation. This section summarises the
conclusions of chapter~\ref{ch:evaluation}, focusing on the findings that hold
regardless of the underlying hardware.

\paragraph{Data movement.} For small k, every combination of algorithm and
architecture is limited by bandwidth and not by the computational capability of
the devices (section~\ref{sec:eval-roofline}). Performance depends solely on
reducing the volume of data moved. Bitonic sort on the graphics unit confirms
this clearly. Although it uses a larger fraction of the bandwidth ceiling than
the reference library, it ultimately falls behind because it is forced to move
several times as much data.

\paragraph{The ratio $n/k$.} The fastest implementation for k=8 becomes the
slowest for k=131072 regardless of platform (section~\ref{sec:eval-time-k}).
The performance of the sorting network stays far more stable. The heap
algorithm starts filtering elements only once it is full, which caps the
available parallelism by the ratio $n/k$. Similar behaviour has also been
observed in the literature~\cite{lin2026adaptivetopk}, highlighting the need
for adaptive algorithm switching depending on the input data.

\paragraph{The measurement boundary.} While the graphics unit wins in pure
kernel execution time at the extreme sizes of the experiment, the overall
end-to-end measurement shows the processor to be faster in every case. Shifting
the measurement boundary completely changes the final ranking. The evaluation
of an implementation must account for the cost of data transfer; otherwise the
comparison of whether offloading to the card pays off remains incomplete.

\paragraph{Power and energy.} The neural processing unit has the lowest power
draw of all implementations and at the same time the highest consumption, since
duration dominates the product (section~\ref{sec:eval-energy}).

For top-k selection with small values of k, offloading the computation to an
accelerator proves disadvantageous in both energy and time compared to running
map-reduce on the processor. Architectures such as NPUs
(section~\ref{sec:bg-arch-npu}) presuppose high local data reuse to amortise
their transfer cost. Top-k selection does not offer this, so most of the time
is spent preparing the data, leaving the device idle
(section~\ref{sec:npu-pipeline}).

The only exception is bitonic sort on the neural processing unit for
half-precision numbers, where it surpasses the corresponding reference
implementation. The advantage is attributed to the incomplete optimisation of
the processor's library for this particular data type
(section~\ref{sec:eval-dtypes}).

\section{Limitations}
\label{sec:concl-limitations}

The comparison with std::partial\_sort on the processor involves an inherent
asymmetry. The library runs strictly single-threaded, unlike the
implementations of this work, which exploit parallelism
(section~\ref{sec:cpu-reference}). The recorded speed-up results jointly from
both the algorithmic improvement and the use of multiple threads. This choice
stems from the intention to compare the new implementation with the tools a
programmer has immediately available, but the difference should be kept in mind
when reading the results.

In addition, the volume of data counted for the reference libraries is merely a
lower bound. Because their internal memory passes are not directly observable,
the ceiling-utilisation percentages reported in
section~\ref{sec:eval-bandwidth} are underestimates.

For the neural processing unit no native top-k library is available for direct
comparison. The comparison was made against the host, answering the practical
question of whether offloading is worthwhile, without evaluating the maximum
theoretical performance of the device itself. Correspondingly, the phase
analysis of section~\ref{sec:npu-pipeline} comes from separate runs with debug
enabled rather than from the standard sweep, so its percentages rest on a
smaller sample of repetitions than the other measurements of
appendix~\ref{app:reproducibility}.

Finally, the experimental space was limited to a specific system and specific
ranges of values. Although the qualitative trends around the ratio $n/k$ and
the role of the measurement boundary are expected to generalise to similar
systems, the exact numerical factors remain closely tied to this particular
test platform.

\section{Future work}
\label{sec:concl-future}

In terms of system-level optimisations, executing multiple simultaneous top-k
queries is the main future extension, especially for the neural processing
unit. Since the main obstacle lies in preparing and dispatching the data,
batching the calls would help decisively in amortising this cost. Such repeated
queries dominate the applications of section~\ref{sec:bg-problem}. In the same
direction, exploiting mechanisms for direct device access to system memory,
avoiding explicit copying of data by the host, could significantly improve
overall latency. Implementing such a capability depends to a large extent on
the specifications and interfaces of each platform.

As for algorithmic extensions, a mechanism for dynamic algorithm switching
would directly exploit the finding about the ratio $n/k$. Knowing now how the
ratio affects performance, an adaptive implementation could automatically
choose between a heap and a sorting network to ensure smooth performance across
the whole range of sizes. Furthermore, for larger values of $k$, where
selection approaches full sorting, handling inputs that do not fit in device
memory is a natural continuation, turning stream management from a
responsibility of the harness into an organic part of the algorithm itself.

Finally, the evaluation harness developed in this work can be used in the
future to study other workloads as well. Separating algorithmic time from
transfer time, recording energy uniformly across the different architectures
and analysing through the roofline model offer a reusable framework, capable of
being applied to any other problem considered for offloading to heterogeneous
hardware.
